\documentclass[eprint,
%approc,		%heading for Acta Polytechnica CTU Proceedings
]{actapoly}

\usepackage[utf8]{inputenc}

\newcommand*\printft[1]{\textsc{\lowercase{#1}}}

\begin{document}

\title[Pareto Efficiency and RSM to Adjust Binder Content]
{Pareto Efficiency and RSM to Adjust Binder Content for Clay Stabilization in Yttre Ringvägen, Malmö}

\author[P. Lindh]{Per Lindh}{my,their}
\correspondingauthor[P. Lemenkova]{Polina Lemenkova}{third}{polina.lemenkova@ulb.be}

\institution{my}{Swedish Transport Administration, Gibraltargatan 7, Malmö, Sweden}
\institution{their}{Lund University, Division of Building Materials, Box 118, SE- 221-00, Lund, Sweden}
\institution{third}{Université Libre de Bruxelles (ULB), École polytechnique de Bruxelles (Brussels Faculty of Engineering), Laboratory of Image Synthesis and Analysis (LISA). Campus de Solbosch, ULB -- LISA CP165/57, Avenue Franklin D. Roosevelt 50, B-1050 Brussels, Belgium}

\begin{abstract}
In this paper, we present a new framework for improvement of engineering properties of soil before road construction utilizing an advanced method of stabilization by blended binders. The effects of various binder combinations and water content on the performance of soil strength were evaluated using statistical methods. Specimens included clay till collected from the Yttre Ringvägen, Sweden. We applied quicklime hydrated lime, cement, slag as pure binders and their blended mixes. Clay tills were tested and the effects from change in water and binder content before-after stabilisation was shows using Pareto charts. The statistical analysis included quadratic model, regression analysis and ANOVA to analyse the efficiency and quality of testing methods. We examined two different variables: 1) the effects from the amount of binder on soil performance; 2) the interaction between several binders as changed ratio of cement/lime/slag: 30-50-20\%; 50-50-0\%; 100-0-0\%. Three Response Surface Methodology (RSM) techniques were evaluated, compared, and presented in this experiment: 1) Central Composite Design (CCD); 2) Box-Behnken Design (BBD); 3) Simplex Lattice Design. The increase of cement improved the stabilization performance followed by slag and lime.

soil stabilisation, simplex experimental design, binder, cement, statistical analysis.

\end{abstract}

\keywords{Keyword one, keyword two, keyword three}

\maketitle


\section{Introduction}

Soil stabilisation is widely applied method of improving soil parameters and performance in civil engineering. It can be performed using various binders as stabilising agents. Despite the existing methods, new stabilisation agents and approaches require experimental testing and evaluation to assess their quality and effectiveness. 

\section{Materials and methods}

\newpage
\section{Discussion and Conclusions}

In this paper, we introduced a framework for selecting optimal amounts of binders used for soil stabilization with a case study of clay soils collected from the region of Yttre Ringvägen, a ring road in Malmö, southern Sweden. A novel framework for adjusting binder content was proposed. Three approaches of the RSM techniques were tested for evaluation of soil performance stabilised using various binder types. The results were compared and presented using Pareto charts. The methods assessed in this study area: 1) Central Composite Design (CCD); 2) Box-Behnken Design (BBD); 3) Simplex Lattice Design. 

We used clay tills from the embankment fill material collected from the test places connecting road for the Öresund bridge, southern Scania region of Sweden. The performed experimental design proved to be a perfect tool for improving the efficiency and quality of various methods of testing soil quality and performance after stabilization with various binders: cement, slag, and lime. Besides, the aforementioned comparisons of various stabilising agents and the performance of soil upon treatment we also conducted the statistically significant test using the ANOVA approach implemented in Statistica software. From the achieved results, we can conclude that the difference between the testing methods is statistically significant (p-values are all smaller than 0.05 significance level). Namely, we noted that cement, lime and the interaction between cement and lime are significant at p = 0.05 and that the quadratic effects of cement and lime are not significant.

To demonstrate the capability of various binder and their influence on soil stabilization and strength development, two approaches were tested: a design for evaluating the amount of binder and a design to evaluate the interaction between several binders. The proposed method is capable of incorporating the RMS which as an effective tool for geotechnical experiments. We showed that the CCD using the robust functionality is essential for assessment of the amount and type of stabilizing agents. Additionally we indicated that BBD can be used for statistical tests. 

Normal factor designs is recommended to be avoided if single stablising agents are selected. This is owing to the fact that zero value of normal factor includes the raw specimens. In case it oversteps zero value, only blended binders should be tested. Therefore, our method can be adopted if the goal of the test is to adjust the most suitable content of binders for soil stabilization, including specimens collected from other regions. We evaluated the mixture design of the ternary statistical modelling  on mixtures of binders for soil stabilization, which proved to be effective and robust approach. Nevertheless, the type of soils and variation in water content in sample specimens should always be taken in the consideration for geotechnical engineering prior to construction works. 

The specific value of our research is that the methods and techniques have been borrowed from other industrial areas such as the automotive, chemical and process industries. Thus, we used the RSM as an advanced and efficient technique which can be used to optimise binder type and amount in soil stabilisation. Using this approach, we have conducted extensive experiments on soil stabilisation using new binder agents and novel methods of soil stabilisation which required different testing and evaluation methods in order to benchmark the performance of the framework. 

\begin{nomenclature}
\item{UCS}{Unconfined Compressive Strength}
\item{BBD}{Box-Behnken Design}
\item{ANOVA}{ANalysis Of VAriance}
\item{CCD}{Central Composite Design}
\item{DMM}{Deep Mixing Method}
\item{ICL}{Initial Consumption of Lime}
\item{MCV}{Moisture Compaction Value}
\item{OMC}{Optimum Moisture Content}
\item{RSM}{Response Surface Methodology}
\item{SGI}{Swedish Geotechnical Institute}
\end{nomenclature}

\begin{acknowledgements}
List funding sources according to the funder’s requirements.
\end{acknowledgements}


\bibliographystyle{actapoly}
\bibliography{biblio}
\nocite*{}

\end{document}
